Three different coordinates system are considered here. Unit vectors are
covered with a hat like $\widehat{v}$. 

\begin{itemize}
\item  First, the toroidal coordinate system $\left( R,Z,\phi \right) $ with
local orthogonal vectors 
\begin{eqnarray}
\widehat{R} &=&\nabla R \\
\widehat{Z} &=&\nabla Z \\
\widehat{\phi } &=&R\nabla \phi 
\end{eqnarray}
This coordinate system conserves the largest generality in the magnetic
geometry

\item  Second, the poloidal (polar) coordinate system $\left( r,\theta ,\phi
\right) $, centered around the magnetic axis $\left( R_{0},Z_{0}\right) $
defined as the position of an extremum of $\psi $ (can be arbitrarily chosen
as $\psi =0$), and with local orthogonal vectors 
\begin{eqnarray}
\widehat{r} &=&\nabla r \\
\widehat{\theta } &=&r\nabla \theta  \\
\widehat{\phi } &=&R\nabla \phi 
\end{eqnarray}
and the relations 
\begin{eqnarray}
R &=&R_{0}+r\cos \theta  \\
Z &=&Z_{0}+r\sin \theta 
\end{eqnarray}
which can be inverted as 
\begin{eqnarray}
r &=&\sqrt{\left( R-R_{0}\right) ^{2}+\left( Z-Z_{0}\right) ^{2}} \\
\theta  &=&\arctan \left( \frac{Z-Z_{0}}{R-R_{0}}\right) +\pi \text{sign}%
\left( R-R_{0}\right)   \label{theta}
\end{eqnarray}

\item  Third, the flux coordinate system $\left( \psi ,s,\phi \right) $ is
the natural system for Fokker-Planck simulations, since the particles are
essentially confined to a given flux surface $\psi $. The distance $s$ is
the length along the poloidal magnetic field lines. We can choose its origin
as being at the position of minimum $B$-field amplitude within a
flux-surface. 
\begin{equation}
B\left( \psi ,s\equiv 0\right) =\min \left\{ B\left( \psi ,s\right) \right\}
=B_{\min }\left( \psi \right) 
\end{equation}
The direction is counter-clockwise and the limits $s_{\min }$ and $s_{\max }$
are set at the position of maximum magnetic field

\begin{equation}
B\left( \psi ,s\equiv s_{\min }\right) =B\left( \psi ,s\equiv s_{\max
}\right) =\max \left\{ B\left( \psi ,s\right) \right\} =B_{\max }\left( \psi
\right)   \label{bmax}
\end{equation}
The local orthogonal vectors are 
\begin{eqnarray}
\widehat{\psi } &=&\frac{\nabla \psi }{\left| \nabla \psi \right| }
\label{psih} \\
\widehat{s} &=&\frac{\nabla \phi \times \nabla \psi }{\left| \nabla \psi
\right| \left| \nabla \phi \right| }=-\widehat{b}_{p} \\
\widehat{\phi } &=&R\nabla \phi   \label{phih}
\end{eqnarray}
where $\widehat{b}_{p}$ is the unit vector in the direction of the poloidal
magnetic field.

\item  An alternative to the previous system may be easier to implement
numerically for the calculation of the bounce coefficients. This is the
system $\left( \psi ,\theta ,\phi \right) $. The $\theta $ grid is defined
in (\ref{theta}). We also define, for geometrical purposes, a flux-function $%
\rho \left( \psi \right) $ which coincides with the normalized radius on the
horizontal LFS midplane. Indeed, in an axisymmetric system, using the
functions $R\left( \psi ,\theta \right) $ and $Z\left( \psi ,\theta \right) $%
, we define $\rho \left( \psi \right) $ as 
\[
\rho \left( \psi \right) =\frac{R\left( \psi ,0\right) -R_{0}}{R_{\max
}-R_{0}}
\]
where $R_{\max }=R\left( \psi _{\max },0\right) $ is the value of $R$ on the
separatrix as it crosses the midplane. The local orthogonal vectors are the
same as (\ref{psih}-\ref{phih}). The 2-D outputs from the axisymmetric
equilibrium code HELENA\ are given on the $\left( \psi ,\theta \right) $
grid. The system $\left( \psi ,\theta ,\phi \right) $ will be used from now
on.
\end{itemize}
